%!TEX root = ../main.tex

\section{Уравнения Максвелла}

Уравения Максвелла в электростатическом приближении можно записать в следующей форме, предварительно исключив из рассмотрения величины, связанные с магнитным полем.
\begin{subequations}
\begin{gather}
    \vE = -\nabla \Phi,
    \label{eq:electrostatic_a} \\
    \partt{\rho} + \Div (\sigma \vE) = 0,
    \label{eq:electrostatic_b} \\
    \Div (\epsilon \vE) = \rho.
    \label{eq:electrostatic_c}
\end{gather}
\label{eq:electrostatic}
\end{subequations}

Вектор Умова--Пойнтинга $\vect{S} = [\vE, \vect{H}]$ есть вектор переноса энергии электромагнитного поля. В электростатическом приближении верно
\begin{equation}
	\Div \vect{S} = \Div [\vE, \vect{H}] = -(\vE, \Rot \vect{H}) = -\left( \vE, \partt{\vD} \right) - (\vE, \vj),
	\label{eq:poynting_vector}
\end{equation}
где слагаемое
\begin{equation}
	(\vE, \vj) = Q_j
	\label{eq:heat}
\end{equation}
выражает джоулево тепло, выделяющееся при протекании тока в среде.